Lo scopo di questa tesi magistrale è di derivare le equazioni effettive per un modello microscopico (un sistema di equazioni differenziali stocastiche) per la chemotassi di $N$ amebe attratte dalla molecola di cAMP (adenosina monofosfato ciclico) che esse stesse secernono in istanti ditribuiti casualmente nel tempo tramite un processo di Poisson.
L'interazione è espressa mediante un potenziale di tipo Gaussiano.
Il modello microscopico è ispirato a \cite{ana}, che, insieme a \cite{peter-boers}, è la fonte principale delle tecniche adottate nella dimostrazione del risultato principale.
Inoltre discutiamo delle somiglianze e delle differenze che la descrizione macroscopica che abbiamo derivato ha con il modello macroscopico classico per la chemotassi, cioè l'equazione di \emph{Keller-Segel}.